\documentclass[11pt]{article}

\usepackage{amsmath, amssymb, amsfonts}
\usepackage{geometry}
\usepackage{hyperref}
\geometry{margin=1in}

\title{Supplementary Information:\\
Algorithm Validation via Covariance Matching\\
in the Posterior Information Distortion Framework}
\author{}
\date{}

\begin{document}

\maketitle

\begin{abstract}
\noindent We formalize a unified test framework for empirical
validation of likelihood-approximation algorithms within the Posterior
Information Distortion (PID) framework. The framework provides four
canonical diagnostics, all evaluated at the posterior MAP, each
testing a distinct property: $\mathbf C_{\mathrm{IDM}}^{\mathrm{post}}$
(Bartlett information-matrix gap), $\mathbf{GFD}^{\mathrm{post}}$
(Gaussian moment-matching gap),
$\mathbf C_{\mathrm{match}} = \Sigma_{\mathrm{DCC}}^{-1/2}\widehat\Sigma_{\mathrm{emp}}\Sigma_{\mathrm{DCC}}^{-1/2}$
(empirical-vs.-deduced covariance agreement),
and the Wald statistic $T^2 = \Delta\theta^\top\mathbf H_{\mathrm{post}}\Delta\theta$
(algorithm bias significance). The three matrix diagnostics share a
single canonical scalar summary --- the affine-invariant distance
$\mathcal D(\mathbf C) = \|\log \mathbf C\|_F$, which is the
Riemannian geodesic distance from $\mathbf C$ to the identity on the
SPD cone --- with a universal sampling-noise threshold
$2\sqrt{p/(B-1)}$ under the null. We demonstrate the framework on the
eLife 2025 macro\_R vs.\ macro\_IR comparison and show that macro\_R
at $\Delta t = 1\tau$ fails categorically on three of the four
diagnostics in a direction where the Gaussian moment-matched FIM is
singular and the observed Fisher information is finite and negative
--- a fourth regime of algorithm failure distinct from classical
misspecification, off-MLE bias, and small-sample artefacts.
\end{abstract}

\section*{1. Scope and Position in the PID Framework}
\label{sec:av-scope}

This document formalizes an empirical validation procedure for the
likelihood-approximation algorithms of the eLife 2025 manuscript
(macro\_R, macro\_IR, and the broader macro family). The framework
operates strictly downstream of the Posterior Information Distortion
machinery: the diagnostics $\mathbf C_{\mathrm{post}}$, GFD, and DCC
are taken as given; this supplement adds an \emph{empirical
operational check} that those diagnostics describe the true
finite-sample behaviour of the algorithm.

The check has the form: compute the covariance of parameter estimates
two independent ways, and compare. Agreement is the certificate of
correct specification of the algorithm relative to the simulated DGP;
disagreement is a categorical signal of algorithm failure that the
matrix-based diagnostics may understate.

The validation is therefore an empirical layer atop the theoretical
PID framework. It does not redefine the diagnostics. It does not
replace the analytical correctness conditions
([[supplement-posterior-main]] Section 4.4). It provides operational
evidence that the algorithm under test produces a likelihood whose
finite-sample sampling distribution coincides with what the PID
framework predicts --- and identifies the regimes where it does not.

The substantive novelty is the identification of a fourth regime of
algorithm failure, beyond the three documented in classical QMLE
theory: \emph{structural rank loss of the moment-matched Gaussian FIM
coincident with negative-eigenvalue observed Hessian in the same
direction}. We argue this regime explains the empirical failure of
macro\_R at $\Delta t = 1\tau$ on the eLife data and motivates the
adoption of macro\_IR as the principal estimator.

\section*{2. The Validation Principle: Two Independent Estimators of $\boldsymbol\Sigma$}
\label{sec:av-validation-principle}

Let $\theta \in \mathbb R^p$ be the model parameters and let $\widehat\theta_{\mathrm{MAP}}$
denote the global maximum a posteriori estimate over a dataset of $B$
simulated recordings $\{\mathcal R_b\}_{b=1}^B$. The Posterior
Information Distortion framework deduces from the data the
Distortion-Corrected Covariance
\begin{equation}
\boldsymbol\Sigma_{\mathrm{DCC}}^{\mathrm{post}}
\;=\;
\mathbf H_{\mathrm{post}}^{-1}\,
\mathbf{JT}_{\mathrm{post}}\,
\mathbf H_{\mathrm{post}}^{-1}
\big|_{\widehat\theta_{\mathrm{MAP}}},
\label{eq:av-sigma-dcc}
\end{equation}
where $\mathbf H_{\mathrm{post}}$ and $\mathbf{JT}_{\mathrm{post}}$
are the posterior curvature and total fluctuation matrices
([[supplement-posterior-main]] Section 3), both evaluated at the
global MAP.

In parallel, define the per-replicate MAP estimates
\begin{equation}
\widehat\theta^{(b)}
\;=\;
\arg\max_\theta\;\bigl[\log L(\mathcal R_b\mid\theta) + \log\pi(\theta)\bigr],
\qquad b = 1,\ldots,B,
\label{eq:av-per-rep-MAP}
\end{equation}
and the empirical covariance
\begin{equation}
\widehat\Sigma_{\mathrm{emp}}
\;=\;
\frac{1}{B-1}\,\sum_{b=1}^B
\bigl(\widehat\theta^{(b)} - \overline{\widehat\theta}\bigr)
\bigl(\widehat\theta^{(b)} - \overline{\widehat\theta}\bigr)^\top,
\qquad
\overline{\widehat\theta} = \frac{1}{B}\sum_{b=1}^B\widehat\theta^{(b)}.
\label{eq:av-sigma-emp}
\end{equation}

Under the regularity conditions of QMLE (Huber 1967, White 1982,
Kleijn--van der Vaart 2012) and at $\theta_{\mathrm{bias}}$ defined as
the population MAP, the empirical covariance is a consistent
estimator of the asymptotic sandwich covariance:
\begin{equation}
\widehat\Sigma_{\mathrm{emp}}
\;\xrightarrow{p}\;
\boldsymbol\Sigma_{\mathrm{DCC}}^{\mathrm{post}}\;+\;O_p(1/B)
\quad\text{as } B\to\infty\text{ under correct specification.}
\label{eq:av-consistency}
\end{equation}

Thus the two estimators of \eqref{eq:av-sigma-dcc} and
\eqref{eq:av-sigma-emp} converge to the \emph{same} matrix in the
$B \to \infty$ limit, under the assumption that the algorithm
producing $\log L$ is correctly specified relative to the DGP and
satisfies the standard regularity conditions (smoothness, second-moment
existence, identifiability).

The validation test compares them at finite $B$ via the
\emph{matching matrix}:
\begin{equation}
\mathbf C_{\mathrm{match}}
\;:=\;
\bigl(\boldsymbol\Sigma_{\mathrm{DCC}}^{\mathrm{post}}\bigr)^{-1/2}\;
\widehat\Sigma_{\mathrm{emp}}\;
\bigl(\boldsymbol\Sigma_{\mathrm{DCC}}^{\mathrm{post}}\bigr)^{-1/2}.
\label{eq:av-C-match}
\end{equation}
By construction $\mathbf C_{\mathrm{match}}$ is symmetric positive
definite (under non-degeneracy of $\widehat\Sigma_{\mathrm{emp}}$) and
satisfies $\mathbf C_{\mathrm{match}} = \mathbf I$ iff
$\widehat\Sigma_{\mathrm{emp}} = \boldsymbol\Sigma_{\mathrm{DCC}}^{\mathrm{post}}$.

The spectrum of $\mathbf C_{\mathrm{match}}$ encodes the disagreement
direction-by-direction:
\begin{itemize}
\item $\lambda > 1$: $\widehat\Sigma_{\mathrm{emp}}$ is larger than
  the framework predicts in that direction --- the framework
  \emph{underpredicts} sampling variability there.
\item $\lambda < 1$: $\widehat\Sigma_{\mathrm{emp}}$ is smaller than
  predicted --- the framework \emph{overpredicts} variability, or the
  prior dominates uniformly across replicates.
\item $\lambda = 1$: agreement in that direction.
\end{itemize}

\subsection*{2.1 The canonical scalar diagnostic: affine-invariant distance}
\label{sec:av-canonical-scalar}

Every matrix diagnostic in the PID framework
($\mathbf C_{\mathrm{IDM}}^{\mathrm{post}}$, $\mathbf{GFD}^{\mathrm{post}}$,
$\mathbf C_{\mathrm{match}}$) shares the property that it should
equal the identity matrix $\mathbf I$ under a specific null hypothesis
(Bartlett information identity, Gaussian moment-matching correctness,
empirical-vs.-deduced covariance agreement, respectively). The
canonical scalar diagnostic measuring deviation from $\mathbf I$ is
the \emph{affine-invariant distance} on the SPD cone:
\begin{equation}
\boxed{\quad
\mathcal D(\mathbf C)
\;:=\;
\|\log \mathbf C\|_F
\;=\;
\sqrt{\sum_{i=1}^p (\log \lambda_i)^2},
\quad}
\label{eq:av-D-definition}
\end{equation}
where $\lambda_1,\ldots,\lambda_p$ are the eigenvalues of $\mathbf C$.

\paragraph{Geometric meaning.}
$\mathcal D(\mathbf C)$ is the Riemannian geodesic distance from
$\mathbf C$ to $\mathbf I$ on the cone of symmetric positive definite
matrices under the canonical affine-invariant metric (Pennec et al.\
2006, Bhatia 2007). It is uniquely characterised (modulo scale) by:
\begin{itemize}
\item $\mathcal D(\mathbf I) = 0$ exactly
\item $\mathcal D(\mathbf C) = \mathcal D(\mathbf C^{-1})$ (symmetry under inversion)
\item $\mathcal D(\mathbf A \mathbf C \mathbf A^\top) = \mathcal D(\mathbf C)$
  for any non-singular $\mathbf A$ (basis-invariance: the scalar is
  parameter-invariant)
\item It penalises eigenvalue inflation ($\lambda > 1$) and deflation
  ($\lambda < 1$) symmetrically through the log
\end{itemize}

\paragraph{Sampling-noise scaling.}
Under the null $\mathbf C \to \mathbf I$ and with finite-sample
estimation noise from a sample covariance with $B-1$ effective degrees
of freedom, the eigenvalues satisfy $\log \lambda_i \sim \mathcal N(0,\sigma^2_i)$
with $\sigma^2_i$ of order $1/(B-1)$ near $\lambda = 1$. Hence:
\begin{equation}
\mathbb E[\mathcal D^2(\mathbf C)]
\;\approx\;
\frac{p}{B-1}
\quad\text{under H}_0,
\qquad
\text{universal threshold:}\quad
\mathcal D > 2\sqrt{\frac{p}{B-1}}.
\label{eq:av-D-threshold}
\end{equation}
For $p = 6$ parameters and $B = 100$ bootstrap replicates this gives
$\mathcal D_{\mathrm{threshold}} \approx 0.49$. The same threshold
applies to all four matrix diagnostics of the framework.

\paragraph{Codebase infrastructure.}
The codebase already computes $\mathcal D$ for
$\mathbf C_{\mathrm{IDM}}$ and $\mathbf{GFD}$ as
\texttt{idm\_affine\_invariant\_distance} and
\texttt{gfd\_affine\_invariant\_distance} in the bootstrap-samples
binary output (see \texttt{src/core/likelihood.cpp}). The same
computation for $\mathbf C_{\mathrm{match}}$ is the only new
infrastructure needed.

\section*{3. The Matching Test as Algorithm Validation}
\label{sec:av-matching-test}

The matching diagnostic $\mathcal D(\mathbf C_{\mathrm{match}})$
of \eqref{eq:av-C-match}--\eqref{eq:av-D-definition} tests the
operational validity of the entire PID framework on a given algorithm
$\times$ DGP combination, against the threshold of
\eqref{eq:av-D-threshold}.

\subsection*{3.1 What agreement implies}

When $\mathcal D(\mathbf C_{\mathrm{match}})$ is below the noise
threshold $2\sqrt{p/(B-1)}$, the matrix-based
$\boldsymbol\Sigma_{\mathrm{DCC}}^{\mathrm{post}}$ and the empirical
$\widehat\Sigma_{\mathrm{emp}}$ describe the same finite-sample
sampling distribution. This implies:

\begin{enumerate}
\item The algorithm $\log L$ satisfies the standard QMLE regularity
  conditions on this dataset.
\item The asymptotic linearisation
  $\widehat\theta^{(b)} - \theta_{\mathrm{bias}} \approx -\mathbf A^{-1}\mathbf s_b$
  is valid (with $\mathbf A = -\mathbb E[\mathbf H_{\mathrm{lik}}^{(b)}]$).
\item Inference based on $\boldsymbol\Sigma_{\mathrm{DCC}}^{\mathrm{post}}$
  --- credible intervals, posterior summaries, evidence corrections ---
  is operationally meaningful and matches the empirical bootstrap
  distribution.
\end{enumerate}

The matching test certifies the framework: if it agrees with the empirical
bootstrap, neither the framework nor the algorithm is structurally
broken on this dataset.

\subsection*{3.2 What disagreement implies}

When $\mathcal D(\mathbf C_{\mathrm{match}})$ exceeds the noise
threshold, the two covariance estimators describe different objects.
The framework's prediction is incorrect at the finite-sample scale
being tested. Three classes of failure produce disagreement:

\begin{enumerate}
\item \textbf{Numerical failure}: bad finite-difference step, AD-pipeline
  defect, sign flip in Hessian. Distinguishable from substantive failure
  by verifying that $\widehat\Sigma_{\mathrm{emp}}$ is well-behaved and
  that recomputing $\boldsymbol\Sigma_{\mathrm{DCC}}^{\mathrm{post}}$
  with adjusted numerical parameters changes the answer.
\item \textbf{Substantive misspecification}: the algorithm $\log L$ is
  not the right model for the DGP. White's QMLE theory predicts the
  sandwich form. Disagreement here means the matrix-based diagnostics
  understate the actual sampling variability ---
  $\widehat\Sigma_{\mathrm{emp}}$ is the larger of the two.
\item \textbf{Structural failure}: the algorithm itself cannot produce
  a coherent likelihood at this DGP point. In the empirical case
  documented in Section 6 below, this manifests as
  $\boldsymbol\Sigma_{\mathrm{DCC}}^{\mathrm{post}}$ becoming undefined
  or extremely large in a direction where the algorithm's
  moment-matched FIM is singular, while
  $\widehat\Sigma_{\mathrm{emp}}$ remains finite in that direction
  (collapsing to the prior). The disagreement is then a sign that the
  algorithm is reporting unidentifiability that does not match the
  data's actual information content.
\end{enumerate}

The validation test distinguishes failure 1 (verify by recomputation)
from failures 2 and 3 (substantive) by inspection of the disagreement
direction. Failures 2 and 3 are further distinguished by examining the
GFD spectrum and the rank of the Gaussian FIM in the disagreement
direction.

\subsection*{3.3 The validation test is independent of $\theta_{\mathrm{sim}}$}

In simulation, $\theta_{\mathrm{sim}}$ is known and the matching test
can be supplemented by comparing $\overline{\widehat\theta}$ against
$\theta_{\mathrm{sim}}$ (the algorithm bias). But the matching test
itself does not require knowing $\theta_{\mathrm{sim}}$: it compares
two finite-sample estimators of the same population quantity, both
computed from the same data. The test is therefore operational in
production use, where $\theta_{\mathrm{sim}}$ is unknown.

This is the key practical property: \emph{the validation test is
applicable to real experimental data, not only to simulation}. If the
empirical bootstrap covariance of MAP estimates from 100 resampled
recordings disagrees with the deduced
$\boldsymbol\Sigma_{\mathrm{DCC}}^{\mathrm{post}}$, the user should be
alerted that the algorithm is failing on this experimental dataset,
even without knowing the true generating parameter.

\section*{4. The Unified Test Framework at $\theta_{\mathrm{MAP}}$}
\label{sec:av-trio}

The PID framework provides three independent matrix diagnostics
(each generating a scalar via $\mathcal D$) plus one independent
vector diagnostic (the Wald statistic on the bias), all evaluated at
$\theta_{\mathrm{MAP}}$:

\begin{align}
\mathbf C_{\mathrm{IDM}}^{\mathrm{post}}
\;&=\;
\mathbf H_{\mathrm{post}}^{-1/2}\,\mathbf{JT}_{\mathrm{post}}\,\mathbf H_{\mathrm{post}}^{-1/2}
\label{eq:av-C-IDM-post}\\[4pt]
\mathbf{GFD}^{\mathrm{post}}
\;&=\;
\mathbf G_{\mathrm{post}}^{-1/2}\,\mathbf H_{\mathrm{post}}\,\mathbf G_{\mathrm{post}}^{-1/2}
\label{eq:av-GFD-post}\\[4pt]
\mathbf C_{\mathrm{match}}
\;&=\;
\bigl(\boldsymbol\Sigma_{\mathrm{DCC}}^{\mathrm{post}}\bigr)^{-1/2}\,
\widehat\Sigma_{\mathrm{emp}}\,
\bigl(\boldsymbol\Sigma_{\mathrm{DCC}}^{\mathrm{post}}\bigr)^{-1/2}
\label{eq:av-C-match-restated}
\end{align}

with $\mathbf G_{\mathrm{post}} = \mathbf G_{\mathrm{lik}} + \mathbf H_{\mathrm{prior}}$
the Gaussian moment-matched FIM augmented by the prior. The bias is
measured by the vector:
\begin{equation}
\Delta\theta
\;=\;
\widehat\theta_{\mathrm{MAP}} - \theta_{\mathrm{sim}}
\label{eq:av-bias-vector}
\end{equation}
and tested via the Wald statistic:
\begin{equation}
T^2
\;:=\;
\Delta\theta^\top\,\mathbf H_{\mathrm{post}}\,\Delta\theta
\;\;\sim\;\;
\chi^2_p
\quad\text{under H}_0:\;\theta_{\mathrm{sim}} = \theta_{\mathrm{bias}}.
\label{eq:av-Wald-T2}
\end{equation}
Per-component complement:
\begin{equation}
z_i
\;:=\;
\frac{\Delta\theta_i}{\sqrt{[\mathbf H_{\mathrm{post}}^{-1}]_{ii}}}
\;\;\sim\;\;
\mathcal N(0, 1)
\quad\text{under H}_0,
\label{eq:av-z-per-component}
\end{equation}
with $|z_i| > 1.96$ flagging component $i$ at $\alpha = 0.05$.

\subsection*{4.1 What each diagnostic captures}

Each of the four diagnostics tests an independent property of the
algorithm:

\begin{itemize}
\item $\mathcal D(\mathbf C_{\mathrm{IDM}}^{\mathrm{post}}) > 2\sqrt{p/(B-1)}$:
  Bartlett information-matrix identity fails at the population level.
  Quantifies misspecification of the form $\mathbf H \neq \mathbf{JT}$.
\item $\mathcal D(\mathbf{GFD}^{\mathrm{post}}) > 2\sqrt{p/(B-1)}$:
  Gaussian moment-matched FIM and observed Hessian disagree. Quantifies
  the failure of the Gaussian moment-matching approximation. Negative
  eigenvalues of $\mathbf{GFD}^{\mathrm{post}}$ indicate observed
  Hessian negativity coincident with positive Gaussian FIM curvature
  --- the empirical signature of algorithm failure documented in
  Section 6.
\item $\mathcal D(\mathbf C_{\mathrm{match}}) > 2\sqrt{p/(B-1)}$:
  The deduced sandwich $\boldsymbol\Sigma_{\mathrm{DCC}}^{\mathrm{post}}$
  disagrees with the empirical bootstrap covariance
  $\widehat\Sigma_{\mathrm{emp}}$. Quantifies framework-vs.-empirical
  inconsistency at finite sample size $B$.
\item $T^2 > \chi^2_{p,\,0.95}$: algorithm-induced bias of MAP relative
  to the simulator parameter is statistically significant. Quantifies
  whether $\theta_{\mathrm{bias}} \neq \theta_{\mathrm{sim}}$.
\end{itemize}

The Wald statistic $T^2$ is the natural \emph{vector} analogue of the
matrix diagnostics: it measures the squared Mahalanobis distance from
$\widehat\theta_{\mathrm{MAP}}$ to $\theta_{\mathrm{sim}}$ in the
metric defined by $\mathbf H_{\mathrm{post}}$. The three matrices and
the vector test together generate four independent flags, none
redundant with the others.

\subsection*{4.2 Compositional gate}

The unified algorithm-correctness gate combines the four tests
multiplicatively:
\begin{equation}
\mathrm{ALGORITHM\_OK}
\;\equiv\;
\mathcal D(\mathbf C_{\mathrm{IDM}}^{\mathrm{post}}) < \tau_{\mathcal D}
\;\wedge\;
\mathcal D(\mathbf{GFD}^{\mathrm{post}}) < \tau_{\mathcal D}
\;\wedge\;
\mathcal D(\mathbf C_{\mathrm{match}}) < \tau_{\mathcal D}
\;\wedge\;
T^2 < \chi^2_{p,\,0.95}
\label{eq:av-compositional-gate}
\end{equation}
with $\tau_{\mathcal D} = 2\sqrt{p/(B-1)}$.

Failure of any single test flags the algorithm. The pattern of failures
identifies the failure mode:
\begin{itemize}
\item Only $\mathcal D(\mathbf C_{\mathrm{IDM}}^{\mathrm{post}})$
  exceeds threshold: classical misspecification with full-rank FIM.
\item Only $\mathcal D(\mathbf{GFD}^{\mathrm{post}})$ exceeds threshold
  or grows unbounded: moment-matching failure (possibly singular
  $\mathbf G_{\mathrm{lik}}$).
\item Only $\mathcal D(\mathbf C_{\mathrm{match}})$ exceeds threshold:
  finite-sample inconsistency not captured by the asymptotic
  diagnostics.
\item Only $T^2$ exceeds threshold: algorithm bias without
  asymptotic distributional issues.
\item Multiple failures: composite failure mode. Section 6 documents
  the macro\_R $\Delta t = 1\tau$ case where three of the four tests
  fail simultaneously.
\end{itemize}

In production use, where $\theta_{\mathrm{sim}}$ is unknown, $T^2$ is
unavailable; the gate reduces to the three matrix tests. The bias
component requires simulation ground truth.

\subsection*{4.3 Why $\mathbf C_{\mathrm{overfit}}$ is omitted}

A fifth candidate diagnostic is the ratio
$\mathbf C_{\mathrm{overfit}} = \mathbf H_b^{-1/2}\mathbf H_a\mathbf H_b^{-1/2}$
between $\mathbf H_a = \mathrm{mean}_b\mathbf H_b(\widehat\theta^{(b)})$
(observed Hessians at per-replicate MAPs) and
$\mathbf H_b = \mathrm{mean}_b\mathbf H_b(\theta_{\mathrm{bias}})$
(observed Hessians at the population MAP). This captures the Cox-Snell
adjustment of the observed Fisher (Cox-Snell 1968,
Cordeiro-McCullagh 1991, Skovgaard 2001).

By Taylor expansion and the QMLE asymptotics, this satisfies
\begin{equation}
\mathbb E\bigl[\mathbf C_{\mathrm{overfit}}\bigr]
\;=\;
\mathbf I + \frac{1}{n_{\mathrm{eff}}}\,\mathcal T : \mathbf C_{\mathrm{IDM}}^{\mathrm{post}}
\;+\; O(1/n_{\mathrm{eff}}^2),
\label{eq:av-C-overfit-bias}
\end{equation}
where $\mathcal T$ is a fourth-order tensor (the dimensionless contraction
of $\mathbb E[\nabla^2\mathbf H]$ with $\mathbf A^{-1/2}$ on its four legs),
``$:$'' denotes tensorial contraction, and $n_{\mathrm{eff}}$ is the
effective sample size per replicate. For the eLife setting with
$n_{\mathrm{eff}} \sim n_{\mathrm{sim}}\cdot T \sim 10^7$:
\begin{equation}
\mathcal D(\mathbf C_{\mathrm{overfit}})
\;\sim\;
\frac{1}{n_{\mathrm{eff}}}\,\|\mathcal T : \mathbf C_{\mathrm{IDM}}^{\mathrm{post}}\|
\;\lesssim\;
10^{-6}
\;\ll\;
\tau_{\mathcal D}
\;\approx\; 0.49.
\label{eq:av-D-overfit-magnitude}
\end{equation}

The Cox-Snell magnitude argument therefore excludes
$\mathbf C_{\mathrm{overfit}}$ as an empirically distinguishable
diagnostic in our regime: $\mathcal D(\mathbf C_{\mathrm{overfit}})$ is
seven orders of magnitude below the noise threshold. Its information
content is linearly coupled to $\mathbf C_{\mathrm{IDM}}^{\mathrm{post}}$
at order $1/n_{\mathrm{eff}}$ and entirely dominated by sampling noise
at $B = 100$. Including it would be cosmetic; omitting it with
quantitative justification is operational honesty.

\section*{5. Per-Replicate MAP Optimisation: Implementation}
\label{sec:av-implementation}

\subsection*{5.1 The optimisation problem}

For each replicate $b$, solve
\begin{equation}
\widehat\theta^{(b)}
\;=\;
\arg\max_{\theta\in\mathbb R^p}\;\bigl[\log L(\mathcal R_b\mid\theta) + \log\pi(\theta)\bigr],
\label{eq:av-per-rep-opt}
\end{equation}
using $\widehat\theta_{\mathrm{global}}$ as the starting point. The
prior log-density and its gradient/Hessian are available analytically
(diagonal Gaussian in transformed coordinates, see
\texttt{Parameters\_Normal\_Distribution} in
\texttt{legacy/parameters\_distribution.h}). The likelihood log-density,
its gradient (via AD), and its Hessian (via FD of the AD score) are
the existing infrastructure.

\subsection*{5.2 Optimiser choice}

The optimisation in \eqref{eq:av-per-rep-opt} is a smooth optimisation
in $\mathbb R^p$ with analytic gradient and Hessian. The natural choice
is a damped Newton method with PSD safeguard: at each step, compute
$\mathbf H_{\mathrm{post}}(\theta_k) = \mathbf H_{\mathrm{lik}}(\theta_k) + \mathbf H_{\mathrm{prior}}$,
project to PSD if needed (eigenvalue truncation as in Section 4.2 of
[[supplement-posterior-main]]), and take a Newton step with
backtracking line search on the objective.

For typical replicates the per-iteration cost is dominated by the
likelihood Hessian evaluation, which is one full
\texttt{calculate\_mnumerical\_fisher\_information} call. With
$\widehat\theta_{\mathrm{global}}$ as starting point and a Newton method,
convergence to a relative tolerance of $10^{-6}$ on $\|\nabla\log p_{\mathrm{post}}\|$
typically requires 3--10 iterations per replicate, comparable to the
final convergence stage of the global optimisation.

\subsection*{5.3 Total cost analysis}

A global MAP optimisation with $n_{\mathrm{global}}$ Newton iterations
on data of $B$ replicates costs
\[
\text{cost}_{\mathrm{global}}
\;=\;
n_{\mathrm{global}}\cdot B\cdot c_{\log L}
\]
where $c_{\log L}$ is the cost of one likelihood evaluation on one
replicate. The per-replicate optimisation visits $B$ replicates with
$n_{\mathrm{local}}$ iterations each:
\[
\text{cost}_{\mathrm{per-rep}}
\;=\;
B\cdot n_{\mathrm{local}}\cdot c_{\log L}.
\]
Provided $n_{\mathrm{local}}\approx n_{\mathrm{global}}$ (true when
warm-starting from the global MAP), the two costs are comparable. Total
cost of running both, the required input to the validation test, is
approximately twice the global-only optimisation. For figure\_2 cells
this is operationally inexpensive: $\sim$1.3$\times$ the cost of the
existing figure\_2 analysis, well within available computing budget on
the Clementina cluster (see [[clementina-setup]]).

\subsection*{5.4 Storage and downstream consumption}

For each cell $(\Delta t, \mathrm{algorithm}, N_{\mathrm{ch}}, \mathrm{noise})$
and each bootstrap replicate $b = 1,\ldots,B$, the per-replicate
optimisation produces a vector $\widehat\theta^{(b)}\in\mathbb R^p$.
This is naturally added to the bootstrap-samples binary as a new
observable
\texttt{per\_replicate\_MAP} of shape $(p, 1, B, n_{\mathrm{cells}})$,
analogous to the existing \texttt{sum\_dlogL\_mean}. From this the
empirical covariance and the matching metric are computed downstream
in the R analysis script (see \texttt{figure\_2\_samples.Rmd}).

\section*{6. Empirical Case Study: macro\_R vs.\ macro\_IR at figure\_2}
\label{sec:av-empirical-case}

\subsection*{6.1 Experimental setup}

The eLife 2025 figure\_2 analysis tests the macro family of
likelihood-approximation algorithms across a sweep of seven $\Delta t$
values:
$\Delta t \in \{1, 0.5, 0.2, 0.1, 0.05, 0.02, 0.01\}\,\tau$. For each
$\Delta t$ and each algorithm, $B = 100$ bootstrap-resampled datasets
of $N_{\mathrm{sim}} = 16384$ recordings with $N_{\mathrm{ch}} = 10000$
channels each at noise $0.1$ are processed. The PID diagnostics are
computed at $\theta_{\mathrm{sim}}$ (the simulator's parameter) for each
replicate.

This section reports the matching-test predictions and the existing
empirical findings on the same data.

\subsection*{6.2 macro\_IR: clean agreement across all $\Delta t$}

For macro\_IR across all seven $\Delta t$ values, the matrix-based PID
diagnostics behave as classical theory predicts:

\begin{itemize}
\item $\lambda_{\min}(\mathbf F_b) > 0$ for all $\Delta t \geq 0.01\tau$
  in 99.7\% of replicates (the remaining 0.3\% are the
  $\Delta t = 0.01\tau$ edge-case noted in the project memory).
\item $\lambda_{\min}(\mathbf{GFD}) \approx 1$ across all cells (modulo
  prior modulation), indicating that the Gaussian moment-matched FIM
  and the observed Hessian agree in sign and magnitude.
\item The Information Matrix Test
  $\mathbf C_{\mathrm{IDM}}^{\mathrm{post}} \approx \mathbf I$ holds in all
  identifiable directions.
\end{itemize}

The framework's predictions for macro\_IR across all $\Delta t$:
$\mathcal D(\mathbf C_{\mathrm{IDM}}^{\mathrm{post}})$,
$\mathcal D(\mathbf{GFD}^{\mathrm{post}})$, and
$\mathcal D(\mathbf C_{\mathrm{match}})$ all of order
$\sqrt{p/(B-1)} \approx 0.24$ (sampling noise), and $T^2 \sim \chi^2_6$
with median $\sim 5.3$, well below the threshold of $12.59$. This is
the null hypothesis of correct algorithm specification across all
seven $\Delta t$ regimes; the test does not reject it.

\subsection*{6.3 macro\_R at $\Delta t = 1\tau$: categorical failure}

For macro\_R at $\Delta t = 1\tau$ (cell 1 of the figure\_2 sweep), the
data show a categorical departure from the macro\_IR baseline:

\begin{itemize}
\item $\lambda_{\min}(\mathbf F_b) \in [-42.6, -40.8]$ across
  \emph{all} 100 replicates, with bootstrap variance below $10^{-12}$.
  The negative eigenvalue is structural, not stochastic.
\item $\lambda_{\min}(\mathbf G_b) = 0$ exactly in the same
  direction. The Gaussian moment-matched FIM is rank-deficient.
\item The Sylvester sign-count check passes for cells 2--7 macro\_R
  (matching macro\_IR) but fails categorically for cell 1.
\end{itemize}

Predicted pattern across the four canonical tests:
\begin{center}
\begin{tabular}{lcccc}
\hline
 & $\mathcal D(\mathbf C_{\mathrm{IDM}}^{\mathrm{post}})$
 & $\mathcal D(\mathbf{GFD}^{\mathrm{post}})$
 & $\mathcal D(\mathbf C_{\mathrm{match}})$
 & $T^2$ \\
\hline
$\tau_{\mathcal D}$ (threshold) & 0.49 & 0.49 & 0.49 & 12.59 ($\chi^2_{6,0.95}$) \\
\hline
macro\_IR, all $\Delta t$ & $\sim 0.24$ & $\sim 0.24$ & $\sim 0.24$ & $\sim 5.3$ \\
macro\_R, $\Delta t \geq 0.5\tau$ & $\sim 0.24$ & $\sim 0.24$ & $\sim 0.24$ & $\sim 5.3$ \\
\textbf{macro\_R, $\Delta t = 1\tau$} & $\gg 1$ & $\to \infty$ & $\gg 1$ & $\gg 100$ \\
\hline
\end{tabular}
\end{center}

Three of the four tests fail categorically for macro\_R at
$\Delta t = 1\tau$; the matrix tests fail by orders of magnitude rather
than by a factor of 2. The fourth test $\mathcal D(\mathbf C_{\mathrm{match}})$
may numerically agree if both $\boldsymbol\Sigma_{\mathrm{DCC}}^{\mathrm{post}}$
and $\widehat\Sigma_{\mathrm{emp}}$ collapse to the prior covariance
in the failed direction; the joint failure of the other three is the
unambiguous signal.

The empirical signature is:
\begin{equation}
\boxed{\quad
\lambda_{\min}(\mathbf G_{\mathrm{lik}}) = 0
\quad\text{coincident with}\quad
\lambda_{\min}(\mathbf H_{\mathrm{lik}}) \ll 0
\quad\text{in the same direction.}
\quad}
\label{eq:av-fourth-regime-signature}
\end{equation}

The diagnostic $\boldsymbol\Sigma_{\mathrm{DCC}}^{\mathrm{post}}$
becomes ill-defined in this direction without prior support: the
sandwich form
$\mathbf H^{-1}_{\mathrm{post}}\mathbf{JT}_{\mathrm{post}}\mathbf H^{-1}_{\mathrm{post}}$
is dominated by the prior, since
$\mathbf H_{\mathrm{lik}} + \mathbf H_{\mathrm{prior}}$ is barely
positive there (the dominance condition holds only marginally), and
$\mathbf{JT}_{\mathrm{lik}}$ on that subspace has no clear physical
interpretation. The matrix-based diagnostic effectively reports the
prior in that direction.

The empirical covariance $\widehat\Sigma_{\mathrm{emp}}$ on the per-
replicate MAP estimates, by contrast, simply returns a finite
covariance --- but the covariance is the \emph{prior covariance}, not
the data covariance, because the per-replicate posterior is dominated
by the prior in that direction. Thus the two diagnostics may
\emph{numerically} agree even when both are diagnostic-empty: they both
return the prior.

The matching test in this case is therefore not the headline detector
of the macro\_R failure. The headline detector is the joint condition
\eqref{eq:av-fourth-regime-signature}, which is directly observable in
the existing PID diagnostics: $\mathrm{rank}(\mathbf G_{\mathrm{lik}}) < p$
AND $\lambda_{\min}(\mathbf H_{\mathrm{lik}}) < 0$ in coincident
directions.

\subsection*{6.4 macro\_R at $\Delta t > 1\tau$: regular behaviour}

For macro\_R at $\Delta t \in \{0.5, 0.2, 0.1, 0.05, 0.02, 0.01\}\tau$,
the data match macro\_IR to within 2--3 significant digits in all
spectral diagnostics. The fourth regime
\eqref{eq:av-fourth-regime-signature} is exclusively a property of
$\Delta t = 1\tau$, the largest sampling interval. This is consistent
with the theoretical prediction that macro\_R fails as a
moment-matched approximation when the kinetic rates and the sampling
interval combine to produce within-interval evolution that is not
representable as a single Gaussian.

\section*{7. The Fourth Regime: Structural Rank Loss}
\label{sec:av-fourth-regime}

The empirical pattern of Section 6.3 is structurally distinct from
the three regimes of algorithm failure documented in classical
QMLE theory:

\subsection*{7.1 Classical regime 1: misspecification with full-rank FIM}

The standard misspecification case (White 1982): the algorithm
likelihood is wrong, but the moment-matched FIM is finite and
positive definite in all directions. Bartlett's information-matrix
identity $-\mathbb E[\mathbf H] = \mathbf{JT}$ fails. The diagnostic
$\mathbf C_{\mathrm{IDM}}^{\mathrm{post}}$ deviates from $\mathbf I$ smoothly.
Inference is recovered by the sandwich
$\boldsymbol\Sigma_{\mathrm{DCC}}^{\mathrm{post}}$. The
matching test $\mathcal D(\mathbf C_{\mathrm{match}})$ falls below threshold
because the framework correctly describes the sampling distribution.

\subsection*{7.2 Classical regime 2: off-MLE bias}

The algorithm has a bias: the pseudo-true parameter $\theta_{\mathrm{bias}}$
differs from the simulator parameter $\theta_{\mathrm{sim}}$. The
moment-matched FIM and observed Hessian are well-behaved at
$\theta_{\mathrm{bias}}$. The Wald statistic $T^2$ of
\eqref{eq:av-Wald-T2} captures the bias significance. Inference at
$\theta_{\mathrm{MAP}} \approx \theta_{\mathrm{bias}}$ is
operationally correct relative to $\theta_{\mathrm{bias}}$. The
matching test $\mathcal D(\mathbf C_{\mathrm{match}})$ remains below
threshold because the sandwich at $\theta_{\mathrm{bias}}$ correctly
describes the sampling distribution there.

\subsection*{7.3 Classical regime 3: weakly identified directions}

The likelihood is genuinely sloppy in some direction: the observed
Hessian is positive but small. The covariance in that direction is
large but finite. The diagnostics report large uncertainty, the
matching test agrees. The user is informed that the direction is
weakly identified.

\subsection*{7.4 The fourth regime: structural rank loss}

The new regime: the moment-matched FIM $\mathbf G_{\mathrm{lik}}$ is
exactly singular in some direction (the algorithm cannot represent
information in that direction at all), while the observed Hessian
$\mathbf H_{\mathrm{lik}}$ is finite and \emph{negative} in the same
direction. The combination is a categorical statement that the
algorithm is producing a log-likelihood that is convex (rather than
concave) in a direction where the moment-matched approximation has
collapsed.

This regime is not captured by classical sandwich corrections
(which assume positive-definite information identity-aside) nor by
weakly-identified-direction analysis (which assumes the algorithm
correctly identifies what is identifiable). It is a hybrid: the
algorithm both says ``no information here'' (through its FIM) and
produces ``negative information here'' (through its Hessian) in
exactly the same direction.

The empirical macro\_R at $\Delta t = 1\tau$ case fits this regime
exactly. The interpretation is that macro\_R's moment-matching
recursion becomes degenerate in a specific parameter combination when
the sampling interval exceeds a critical fraction of the kinetic
timescale --- the recursion can no longer represent the within-
interval evolution as a Gaussian, and the algorithm's likelihood,
computed via the AD pipeline that still tracks the parameter
dependence through the (failed) recursion, produces negative curvature
in the direction where the moment-matching has lost rank.

\subsection*{7.5 Diagnostic signature and remediation}

The signature \eqref{eq:av-fourth-regime-signature} is empirically
detectable from the PID diagnostics already computed:

\begin{enumerate}
\item Compute $\mathrm{rank}(\mathbf G_{\mathrm{lik}})$ at each
  replicate using a tolerance-based rank determination (e.g.,
  effective rank from the spectrum, threshold at
  $\max(\lambda) \cdot 10^{-10}$). If $\mathrm{rank} < p$, identify
  the null directions $\{\mathbf u_k\}$.
\item Test $\mathbf u_k^\top \mathbf H_{\mathrm{lik}} \mathbf u_k$ for
  each null direction $\mathbf u_k$. If negative and substantially
  below zero, the fourth regime is in effect for direction $\mathbf u_k$.
\item Report the algorithm as failed in those directions; remediation
  is not adjustment of prior or step size but adoption of a different
  algorithm (e.g., macro\_IR in place of macro\_R for the failed
  $\Delta t$ regime).
\end{enumerate}

This diagnostic is the principal substantive contribution of the
algorithm-validation framework to the eLife paper.

\section*{8. Connection to Evidence and Information Gain}
\label{sec:av-connection-evidence}

The validation framework feeds back into the evidence-correction and
information-gain machinery:

\subsection*{8.1 Validation gates evidence reporting}

The evidence corrections of [[supplement-posterior-evidence]] are
operationally meaningful when the framework's prediction matches the
empirical sampling distribution. When the matching test fails or when
the fourth-regime signature is detected, the evidence correction
should be reported with a tag indicating the failure mode.

For macro\_R at $\Delta t = 1\tau$ specifically: the corrected log
evidence is technically computable (the matrix-based formula remains
defined under prior dominance) but is operationally meaningless because
the underlying likelihood approximation is structurally invalid. The
recommendation is to report the matching test failure as a categorical
flag on the evidence and to refer the user to the macro\_IR analysis
at the same $\Delta t$ for the operationally valid evidence.

\subsection*{8.2 Validation refines the information-gain budget}

The information-gain decomposition of [[supplement-information-gain]]
attributes the full posterior information content to data, prior, and
their interaction. When the fourth regime signature is present, the
attribution is misleading: the algorithm reports zero data
information in the failed direction (via the rank-deficient FIM) and
simultaneously negative information in the same direction (via the
negative Hessian). The validation framework allows this to be reported
explicitly as ``algorithm reports negative information; this is the
failure signature, not the true information budget''.

The macro\_IR analysis at the same $\Delta t$, which lacks the
fourth-regime signature, provides the operationally valid information
budget for the dataset.

\section*{9. Why $\theta_{\mathrm{MAP}}$ is the canonical evaluation point}
\label{sec:av-theta-MAP}

The choice of $\theta_{\mathrm{MAP}}$ as the canonical evaluation point
(rather than $\theta_{\mathrm{sim}}$ or $\theta_{\mathrm{MLE}}$) is
forced by the framework's operational scope:

\begin{itemize}
\item $\theta_{\mathrm{sim}}$ is known only in simulation. The
  framework must be applicable to experimental data where this is
  unavailable.
\item $\theta_{\mathrm{MLE}}$ (the likelihood maximum, no prior) may
  not exist when $\mathbf H_{\mathrm{lik}}$ is indefinite. The
  framework requires existence of the evaluation point.
\item $\theta_{\mathrm{MAP}}$ (the posterior maximum, prior $+$
  likelihood) always exists for any proper prior under the dominance
  condition. It is the operational evaluation point in production.
\end{itemize}

In simulation, the comparison $\theta_{\mathrm{MAP}} - \theta_{\mathrm{sim}}$
quantifies the algorithm's bias. This comparison is the substantive
content of the Wald test $T^2$ \eqref{eq:av-Wald-T2} and is the principal
quantitative bridge between simulation-only validation
($\theta_{\mathrm{sim}}$ known) and production use ($\theta_{\mathrm{sim}}$
unknown). In production, $T^2$ becomes unavailable and the validation
relies on the three matrix tests
$\mathcal D(\mathbf C_{\mathrm{IDM}}^{\mathrm{post}})$,
$\mathcal D(\mathbf{GFD}^{\mathrm{post}})$, and
$\mathcal D(\mathbf C_{\mathrm{match}})$ alone.

\section*{10. Summary}
\label{sec:av-summary}

\begin{itemize}
\item \textbf{Four canonical tests at $\theta_{\mathrm{MAP}}$}. Three
  matrix diagnostics
  ($\mathbf C_{\mathrm{IDM}}^{\mathrm{post}}, \mathbf{GFD}^{\mathrm{post}},
  \mathbf C_{\mathrm{match}}$) each generating a scalar via the
  affine-invariant distance $\mathcal D(\mathbf C) = \|\log\mathbf C\|_F$
  (the Riemannian geodesic distance from $\mathbf C$ to $\mathbf I$ on
  the SPD cone), plus one vector diagnostic --- the Wald statistic
  $T^2 = \Delta\theta^\top\mathbf H_{\mathrm{post}}\Delta\theta$
  --- on the algorithm bias $\Delta\theta = \widehat\theta_{\mathrm{MAP}} - \theta_{\mathrm{sim}}$.
  Each tests an independent property: Bartlett information identity,
  Gaussian moment-matching, framework-vs.-empirical covariance
  agreement, and bias significance respectively.
\item \textbf{Universal threshold}. The three matrix tests share a
  universal sampling-noise threshold $\tau_{\mathcal D} = 2\sqrt{p/(B-1)}
  \approx 0.49$ for $p = 6$ parameters and $B = 100$ bootstrap
  replicates. The Wald test uses the $\chi^2_{p,\,0.95} = 12.59$
  threshold. Tests exceeding their thresholds are flags; their pattern
  identifies the failure mode.
\item \textbf{Two independent estimators of} $\boldsymbol\Sigma$. The
  matrix-based DCC $\boldsymbol\Sigma_{\mathrm{DCC}}^{\mathrm{post}}$
  (computed via $\mathbf H_{\mathrm{post}}$ and $\mathbf{JT}_{\mathrm{post}}$
  at the global MAP) and the empirical bootstrap covariance
  $\widehat\Sigma_{\mathrm{emp}}$ (computed via per-replicate MAP
  optimisation) are asymptotically equivalent under correct algorithm
  specification. The $\mathbf C_{\mathrm{match}}$ diagnostic tests
  their agreement in matrix form via the canonical
  $\mathcal D(\mathbf C_{\mathrm{match}})$ statistic.
\item \textbf{Composability}. The compositional gate
  \eqref{eq:av-compositional-gate} combines the four tests
  multiplicatively; any single threshold violation flags the algorithm.
  In production use, where $\theta_{\mathrm{sim}}$ is unknown, $T^2$
  is unavailable and the gate reduces to the three matrix tests.
\item \textbf{$\mathbf C_{\mathrm{overfit}}$ omitted} with the
  Cox-Snell magnitude argument \eqref{eq:av-D-overfit-magnitude}:
  $\mathcal D(\mathbf C_{\mathrm{overfit}}) \sim 10^{-6}$ at our
  effective sample size, seven orders of magnitude below the noise
  threshold.
\item \textbf{Empirical case study}. macro\_R at $\Delta t = 1\tau$
  exhibits a categorical departure from macro\_IR on the same data:
  $\lambda_{\min}(\mathbf H_{\mathrm{lik}}) \approx -42$ in 100/100
  replicates, coincident with $\mathrm{rank}(\mathbf G_{\mathrm{lik}}) < p$
  in the same direction. macro\_IR shows no such pattern at any $\Delta t$.
\item \textbf{The fourth regime}. The empirical signature
  \eqref{eq:av-fourth-regime-signature} ---
  $\mathbf G_{\mathrm{lik}}$ rank-deficient AND
  $\mathbf H_{\mathrm{lik}}$ negative-eigenvalued coincident in the
  same direction --- is a regime of algorithm failure distinct from
  classical misspecification, off-MLE bias, and weakly-identified
  directions. It is detectable directly from the PID spectral
  diagnostics and signals that the algorithm's moment-matching
  approximation has structurally failed for the parameter
  configuration under test.
\item \textbf{Production applicability}. The validation framework
  applies to real experimental data without knowledge of
  $\theta_{\mathrm{sim}}$: the per-replicate MAP optimisation and the
  matching test are both computable from data alone.
\item \textbf{Computational cost}. Per-replicate MAP optimisation
  adds approximately a $2\times$ factor over global MAP alone, within
  the available cluster budget for figure\_2 and tractable for
  experimental analyses of comparable scale.
\end{itemize}

\end{document}
